World-model-augmented robot policies condition action generation on predicted
future observations or latent visual states.  Pixel prediction is costly and
includes control-irrelevant details, whereas precomputed targets can become
inconsistent with a VLA encoder during fine-tuning.
We introduce \textbf{MINT-VLA} (\textbf{M}ilestone \textbf{I}ntegration with
\textbf{N}ative-Space \textbf{T}argets for
\textbf{V}ision--\textbf{L}anguage--\textbf{A}ction models), a conditioning
adapter that predicts a future task-stage representation in a VLA's current
visual feature space.  MINT-VLA identifies recurring task stages in successful
demonstrations.  During training, the policy encoder jointly encodes the
observation and a future milestone image, and a compact predictor
supplies the estimated milestone through the VLA's native conditioning input.
We instantiate this design as one prefix token in $\pi_{0.5}$.  At inference,
it adds 0.25\% parameters without retrieval, a second visual encoder, or an
image decoder.  In a RoboTwin pilot, MINT-VLA improved a weak no-hint policy by
14.08 percentage points, but exceeded offline latent conditioning by only 0.83
points.  At a fixed offline-conditioning checkpoint, removing the token reduced
success from 70.67\% to 67.17\%, whereas current-state, feature-permuted, and
cross-task replacements reached 71.00\%, 71.50\%, and 71.67\%.  Thus the visual
token affects control, but current evidence does not show that predicted future
content causes the improvement.  MINT-VLA separates future-state prediction
from architecture-specific policy conditioning.
\exptodo{T0--T2,T6,E1}{Replace pilot evidence with the corrected multi-seed
matrix, second-VLA evaluation, and completed matched-scene content
interventions.}
